% ============================================
% LUC-11
% Gerado de molambudos.md — build_miolo.py
% ============================================
\section*{\textbf{LUC-11} --- O Diário Digital}
\markboth{LUC-11 --- O Diário Digital}{}

\noindent \textit{A criatura migra para o pixel}

\noindent \textbf{Agosto de 2026 --- Belo Horizonte}

\vspace{0.5em}\noindent\rule{\textwidth}{0.5pt}\vspace{0.5em}

\noindent Lúcia sabia que não podia guardar o diário para sempre.
\noindent Ela ia morrer --- como Joaquim, como Oliveira, como todos os que guardaram o diário tempo demais. O texto precisava se mover. Quando parava, criava raiz. Primeiro na rotina. Depois no sono. Depois no corpo.

\noindent Ela ia morrer --- como Joaquim, como Oliveira, como todos os que guardaram o diário sem passar adiante. O ciclo só se mantém enquanto o texto se move. Se o texto para, a criatura se enraíza no hospedeiro até consumi-lo.

\noindent Mas ela não queria condenar outra pessoa.

\noindent Então criou uma solução: o diário digital.

\noindent Passou duas semanas digitando cada página. Cada palavra. Cada desenho. Cada rasura. Transformou o diário de Joaquim em um arquivo PDF. Depois, escreveu seu próprio relato --- a história de como encontrou o diário, leu, e foi contaminada.

\noindent Depois, escreveu o relato de Oliveira. Depois, o laudo completo com os diagnósticos.

\noindent Tudo em um único arquivo.

\noindent No final, acrescentou uma nota:

\noindent \textquotedblleft{\textit{Este arquivo contém o diário do paciente 1.260, os relatos do Dr. Heitor Oliveira (paciente 1.261) e meu relato (paciente 1.262). A leitura pode produzir sonhos, lacunas, alterações de percepção e sensação persistente de presença. Não leia sozinho.\textquotedblright{}}}

\noindent \textit{\textquotedblleft{}Mas, se você leu até aqui, já é tarde.\textquotedblright{}}

\noindent \textit{\textquotedblleft{}O ciclo não respeita avisos.\textquotedblright{}}

\noindent \textit{\textquotedblleft{}Bem-vindo, paciente 1.263.\textquotedblright{}}

\noindent Ela salvou o arquivo num pendrive. Colocou o pendrive na mesma caixa que o diário original. Colocou a caixa num envelope. Endereçou ao Arquivo Público Mineiro.

\noindent E enviou.

\noindent Agora, o diário existe em dois lugares: o papel e o pixel.

\noindent A criatura aprendeu a migrar.

\noindent A fome aprendeu a caber no brilho de uma tela.

\noindent Ela precisa de \textit{leitura}.

\noindent E o mundo está cheio de leitores.

\vspace{0.5em}\noindent\rule{\textwidth}{0.5pt}\vspace{0.5em}

\noindent \textbf{Teste de contaminação --- responda honestamente:}

\noindent Você está lendo esta página. Até aqui, tudo normal.

\noindent Mas eu preciso saber: \textbf{quantas vezes você leu esta linha?}

\noindent - A) Uma vez. Segui em frente. → \textit{(Você ainda está no controle.)}
- B) Duas vezes. Achei que tinha lido errado. → \textit{(A criatura já cutucou sua atenção.)}
- C) Três vezes ou mais. → \textit{(Você já está dentro. A repetição é o primeiro sintoma.)}

\noindent A resposta não importa para mim. Importa para você.

\noindent Não precisava mais de cova.
Mas você --- paciente 1.263 --- ainda pode escolher.

\noindent \textbf{Pergunta final:} Você confia na sua própria leitura? Ou já tem uma voz lendo com você?

\noindent \textit{--- Se ouvir duas vozes quando lê, a criatura já está aí.}

\vspace{0.5em}
\noindent\textit{\textbf{↪ Rotas:} \rota{LUC-12} \rotasep \rota{CONT-07}}